\section{De un mal BP a la cadena de suministro del robot completo}

Ahora entramos en el robot de Xinghaitu (星海图). Las secciones anteriores explicaron por qué Gao Jiyang (高继扬) cree en la physical world AI y qué aprendió de la conducción autónoma; esta sección responde a la decisión más importante que tomó al fundar la empresa: por qué una empresa de inteligencia corporizada no puede dedicarse solo a los algoritmos.

El primer giro de Xinghaitu surgió de una conclusión: si se quiere hacer inteligencia corporizada, no basta con hacer la «inteligencia»; hay que fabricar el robot completo. Gao Jiyang dijo en la entrevista que, para que un robot cierre de verdad el ciclo entre datos y modelo, la empresa necesita su propio cuerpo robótico, su hardware, su cadena de suministro y sus escenarios de cliente. De lo contrario, por bueno que sea el modelo base, le falta una vía de acceso estable al mundo real.

\begin{figure}[H]
\centering
\includegraphics[width=0.96\textwidth]{three-product-combo.png}
\caption{Los tres componentes de Xinghaitu: robot completo, modelo base y herramientas de posentrenamiento.}
\end{figure}

\begin{knowledgebox}{Lectura de la figura: por qué tres componentes}
Los tres componentes de la figura apuntan a una misma experiencia objetivo: capacitar a un robot como se capacita a un empleado. El robot completo aporta la vía de acceso física y la capacidad de recopilar datos; el modelo base aporta la base general de acciones y de comprensión; las herramientas de posentrenamiento permiten que el cliente adapte el robot a su escenario con pocas demostraciones y con práctica autónoma. Si falta cualquiera de los tres, es muy difícil llevar el robot de la demo a escenarios productivos.
\end{knowledgebox}

\subsection{El ritmo de 2024, 2025 y 2026}

Los tres componentes describen la estructura final, pero una startup no puede desplegar todas sus capacidades a la vez. Esta subsección se centra en el ritmo, porque el ritmo determina si la organización puede soportar la complejidad.

El ritmo que plantea Gao Jiyang es el siguiente: en 2024 el foco está en el robot completo y la cadena de suministro; en 2025, en la inteligencia de datos; en 2026 empieza el trabajo en la cadena de suministro de largo plazo. Este ritmo refleja la idea de «avanzar paso a paso, consolidando cada posición»: se mantiene la estrategia de largo plazo, pero no se despliega todo al mismo tiempo en un mismo periodo. Si una empresa de robótica persigue demasiado pronto y en paralelo el hardware, los modelos, las aplicaciones, los canales y los clientes, la complejidad termina por desbordar a la organización.

\begin{importantbox}{Implicaciones estratégicas de avanzar paso a paso}
Avanzar paso a paso no es ser conservador, sino avanzar según las dependencias. Sin el robot completo y la cadena de suministro no hay datos estables; sin datos estables no hay un buen VLA; sin un modelo utilizable es difícil generar valor para el cliente.
\end{importantbox}

\subsection{Resumen de la sección}

Xinghaitu eligió fabricar el robot completo no porque el hardware en sí sea más atractivo, sino porque el robot completo es la vía de acceso al ciclo cerrado entre datos, modelo y cliente. Para que una empresa de inteligencia corporizada lleve su inteligencia a fondo, primero tiene que consolidar la cadena física.
